\documentclass[14pt,aspectratio=169,xcolor=dvipsnames]{beamer}
\usetheme{SimplePlus}
\usepackage{booktabs}
\usepackage{minted}
\usepackage{mathtools} % \coloneqq

\title[short title]{Clase 4: Módulos y Python Científico}
\subtitle{}
\author[NA Barnafi] {Nicolás Alejandro Barnafi Wittwer}
\institute[UC|CMM] 
{
    Pontificia Universidad Católica de Chile \\
    Centro de Modelamiento Matemático
}

\titlegraphic{
    \vspace{-1.8cm}
    \begin{flushright}
      \includegraphics[height=2.5cm]{../images/logos/puc.png} 
    \end{flushright}
}

\date{}

\begin{document}
%%%%%%%%%%%%%%%%%%%%%%%%%%%%%%%%%%%%%%%%%%%%%%%%%%%%%%%
\begin{frame}
    \maketitle
\end{frame}
%%%%%%%%%%%%%%%%%%%%%%%%%%%%%%%%%%%%%%%%%%%%%%%%%%%%%%%
%%%%%%%%%%%%%%%%%%%%%%%%%%%%%%%%%%%%%%%%%%%%%%%%%%%%%%%
\section{Módulos}
%%%%%%%%%%%%%%%%%%%%%%%%%%%%%%%%%%%%%%%%%%%%%%%%%%%%%%%
%%%%%%%%%%%%%%%%%%%%%%%%%%%%%%%%%%%%%%%%%%%%%%%%%%%%%%%
\begin{frame}[fragile]\frametitle{Motivación}
    Ya conocemos módulos: \code{string}, \code{sys}
    \begin{itemize}
        \item \code{import string}
        \item \code{import sys}
    \end{itemize}
    Así se cargan ciertas funcionalidades. 

    \pause Hoy aprenderemos a crear nuevos y conoceremos algunos importantes en ciencia
\end{frame}
%%%%%%%%%%%%%%%%%%%%%%%%%%%%%%%%%%%%%%%%%%%%%%%%%%%%%%%
\begin{frame}[fragile]\frametitle{Pip}
    \textbf{P}ackage \textbf{I}nstaller for \textbf{P}ython
    
    \begin{itemize}
        \item Gestor de paquetes que viene instalado junto a Python
        \item Info online en \code{https://pypi.org/}
        \item Fácil de usar por Bash
            \begin{minted}{bash}
        $ pip install pandas
        $ python
        > import pandas  # recién instalada!
            \end{minted}
    \end{itemize}

    \pause Dónde están los programas...?
\end{frame}
%%%%%%%%%%%%%%%%%%%%%%%%%%%%%%%%%%%%%%%%%%%%%%%%%%%%%%%
\begin{frame}\frametitle{el PATH}
    Cuando escribo \code{ls} en Bash... de dónde sale? \\

    \vspace{0.5cm}
    La variable de entorno PATH tiene las carpetas donde buscar programas\footnote{Se puede modificar para tener programas en nuevas carpetas, i.e. \code{/opt}}

    \pause \idea{Consola...( echo \$\{PATH\} )}
\end{frame}
%%%%%%%%%%%%%%%%%%%%%%%%%%%%%%%%%%%%%%%%%%%%%%%%%%%%%%%
\begin{frame}[fragile]\frametitle{Crear módulo}
    \code{operaciones.py:}
    \begin{minted}{python}
  def suma(x,y): return x+y

  def mult(x,y): return x*y
    \end{minted}

  En Python:
  \begin{minted}{python}
  > import operaciones # Debe estar en el path
  > # from operaciones import suma,prod
  > operaciones.suma(2,3) # 5
  > operaciones.mult(2,3) # 6
  \end{minted}
\end{frame}
%%%%%%%%%%%%%%%%%%%%%%%%%%%%%%%%%%%%%%%%%%%%%%%%%%%%%%%
\begin{frame}[fragile]\frametitle{el PATH en Python}
    \begin{minted}{bash}
  ~/clases/Bio210B$ python -m site
   sys.path = [
       '/home/nico/clases/Bio210B',
       '/usr/lib/python3/dist-packages' # son más...
    \end{minted}

    \pause Se puede extender con PYTHONPATH
    \begin{minted}{bash}
  $ export PYTHONPATH=/home/nico/otraCarpeta
  $ python -m site
   sys.path = [
       '/home/nico/clases/Bio210B',
       '/home/nico/otraCarpeta',
       '/usr/lib/python3/dist-packages'
    \end{minted}

\end{frame}
%%%%%%%%%%%%%%%%%%%%%%%%%%%%%%%%%%%%%%%%%%%%%%%%%%%%%%%
\begin{frame}[fragile]\frametitle{Nota: variables en Bash}
    \begin{minted}{bash}
  $ export PYTHONPATH=[...]
    \end{minted}
    genera una variable de ambiente. Se puede ver con
    \begin{minted}{bash}
  $ echo ${PYTHONPATH}
    \end{minted}
  Todas las definidas se ven con programa \code{env} de Bash

\end{frame}
%%%%%%%%%%%%%%%%%%%%%%%%%%%%%%%%%%%%%%%%%%%%%%%%%%%%%%%
%%%%%%%%%%%%%%%%%%%%%%%%%%%%%%%%%%%%%%%%%%%%%%%%%%%%%%%
\section{Python científico}
%%%%%%%%%%%%%%%%%%%%%%%%%%%%%%%%%%%%%%%%%%%%%%%%%%%%%%%
%%%%%%%%%%%%%%%%%%%%%%%%%%%%%%%%%%%%%%%%%%%%%%%%%%%%%%%
\begin{frame}[fragile]\frametitle{Motivación}
    \begin{itemize}
        \item Python es interpretado $\to$ lento
        \item Pero Python usa librerías \emph{MUY rápidas}
        \item Eso + sintaxis lo vuelven ideal para la ciencia
    \end{itemize}
    \begin{center}
    \includegraphics[width=0.35\textwidth]{../images/baby-yoda-T.jpg} \hspace{1cm}\includegraphics[width=0.35\textwidth]{../images/baby-yoda-mesh.png}
    \end{center}
\end{frame}
%%%%%%%%%%%%%%%%%%%%%%%%%%%%%%%%%%%%%%%%%%%%%%%%%%%%%%%
\begin{frame}[fragile]\frametitle{\texttt{math}}
    \begin{minted}{python}
  > import math
  > math. # Tab
math.acos(   math.copysign(  math.expm1(      math.lgamma(     math.pi          math.tanh(
math.acosh(  math.cos(       math.fabs(       math.log(        math.pow(        math.tau
math.asin(   math.cosh(      math.factorial(  math.log10(      math.prod(       math.trunc(
math.asinh(  math.degrees(   math.floor(      math.log1p(      math.radians(    math.ulp(
math.atan(   math.dist(      math.fmod(       math.log2(       math.remainder(  
math.atan2(  math.e          math.frexp(      math.modf(       math.sin(        
math.atanh(  math.erf(       math.fsum(       math.nan         math.sinh(       
math.ceil(   math.erfc(      math.gamma(      math.nextafter(  math.sqrt(       
math.comb(   math.exp(       math.gcd(        math.perm(       math.tan(  
    \end{minted}
    
\end{frame}
%%%%%%%%%%%%%%%%%%%%%%%%%%%%%%%%%%%%%%%%%%%%%%%%%%%%%%%
\begin{frame}[fragile]\frametitle{\texttt{random}}
    \begin{minted}{python}
  > import random
  > random. # Tab
random.BPF           random.SystemRandom( random.gammavariate(   random.paretovariate( random.seed(           
random.LOG4          random.TWOPI         random.gauss(          random.randbytes(     random.setstate(         
random.NV_MAGICCONST random.betavariate(  random.getrandbits(    random.randint(       random.shuffle(          
random.RECIP_BPF     random.choice(       random.getstate()      random.random()       random.triangular(       
random.Random(       random.choices(      random.lognormvariate( random.randrange(     random.uniform(          
random.SG_MAGICCONST random.expovariate(  random.normalvariate(  random.sample(        random.vonmisesvariate( 
    \end{minted}
    
\end{frame}
%%%%%%%%%%%%%%%%%%%%%%%%%%%%%%%%%%%%%%%%%%%%%%%%%%%%%%%
\begin{frame}[fragile]\frametitle{\texttt{statistics}}
    \begin{minted}{python}
  > import statistics as st
  > st. # Tab
st.Counter(          st.bisect_left(  st.fabs(           st.hypot(              st.median(         st.namedtuple( st.repeat(
st.Decimal(          st.bisect_right( st.fmean(          st.itemgetter(         st.median_grouped( st.numbers     st.sqrt(
st.Fraction(         st.correlation(  st.fsum(           st.linear_regression(  st.median_high(    st.pstdev(     st.stdev(
st.LinearRegression( st.covariance(   st.geometric_mean( st.log(                st.median_low(     st.pvariance(  st.tau
st.NormalDist(       st.erf(          st.groupby(        st.math                st.mode(           st.quantiles(  st.variance(
st.StatisticsError(  st.exp(          st.harmonic_mean(  st.mean(               st.multimode(      st.random      
    \end{minted}
    Compatible con arreglos de \code{numpy} 
\end{frame}
%%%%%%%%%%%%%%%%%%%%%%%%%%%%%%%%%%%%%%%%%%%%%%%%%%%%%%%
\begin{frame}\frametitle{NumPy}
    \textbf{Num}erical \textbf{Py}thon
    
    \begin{itemize}
        \item Permite hacer operaciones a grandes datos
        \item Operaciones \emph{vectorizadas}
        \item Submódulos útiles: 
            \begin{itemize}
                \item numpy.fft (Fourier transform)
                \item numpy.linalg (Linear algebra)
                \item numpy.polynomial (Polinomios)
                \item numpy.random (Generador de números)
            \end{itemize}
    \end{itemize}
\end{frame}
%%%%%%%%%%%%%%%%%%%%%%%%%%%%%%%%%%%%%%%%%%%%%%%%%%%%%%%
\begin{frame}[fragile]\frametitle{Numpy arrays}
    \begin{minted}{python}
  > import numpy as np
  > # 50 números equiespaciados en [0,1]
  > num = np.linspace(0,1,50)
  > # Números en [0,1) con paso 0.1
  > num = np.arange(0,1,0.1)
    \end{minted}

\pause \idea{Veamos tiempos de ejecución en sumar}
\end{frame}
%%%%%%%%%%%%%%%%%%%%%%%%%%%%%%%%%%%%%%%%%%%%%%%%%%%%%%%
\begin{frame}[fragile]\frametitle{Operaciones vectorizadas}
    $$ \text{vector:} \qquad x\in \mathbb{R}^N \to \left.\begin{bmatrix} 1\\1.1\\\vdots \\ -0.1\end{bmatrix}\right\}\text{N números} $$

    \begin{minted}{python}
  > import numpy as np
  > num = np.linspace(0,1,50)
  > np.abs(num) # Val absoluto de todo
  > np.exp(num) # exponencial
  > np.sin(num) # seno, ... etc
  > num.dot(num) # normas, producto interno, etc
    \end{minted}

\end{frame}
%%%%%%%%%%%%%%%%%%%%%%%%%%%%%%%%%%%%%%%%%%%%%%%%%%%%%%%
\begin{frame}\frametitle{Matplotlib}
    \centering
    \only<1>{
    \includegraphics[width=0.6\textwidth]{../images/matplotlib/faces.png}
    }
    \only<2>{
    \includegraphics[width=0.5\textwidth]{../images/matplotlib/mandelbrot.png}
    }
    \only<3>{
    \includegraphics[width=0.5\textwidth]{../images/matplotlib/polygon.png}
    }
    \only<4>{
    \includegraphics[width=0.5\textwidth]{../images/matplotlib/time-series.png}
    }
    \only<5>{
    \includegraphics[width=0.5\textwidth]{../images/matplotlib/tricontourf.png}
    }
    \only<6>{
    \includegraphics[width=0.5\textwidth]{../images/matplotlib/3dstem.png}
    }
    \only<7>{
    \includegraphics[width=0.5\textwidth]{../images/matplotlib/3dvoxel.png}
    }
\end{frame}
%%%%%%%%%%%%%%%%%%%%%%%%%%%%%%%%%%%%%%%%%%%%%%%%%%%%%%%
\begin{frame}[fragile]\frametitle{Ejemplo simple}
    \begin{minted}{python}
  > import numpy as np
  > import matplotlib.pyplot as plt
  > xs = np.linspace(0, 2*np.pi, 1000) # [0,2 pi] vectorizado
  > ys = np.sin(xs)
  > plt.plot(xs,ys) # Lista de puntos x e y
  > plt.show()
    \end{minted} 

    \pause Se puede agregar de todo... 

    Ver documentación online: \texttt{matplotlib.org}.
\end{frame}
%%%%%%%%%%%%%%%%%%%%%%%%%%%%%%%%%%%%%%%%%%%%%%%%%%%%%%%
%%%%%%%%%%%%%%%%%%%%%%%%%%%%%%%%%%%%%%%%%%%%%%%%%%%%%%%
\section{Scipy}
%%%%%%%%%%%%%%%%%%%%%%%%%%%%%%%%%%%%%%%%%%%%%%%%%%%%%%%
%%%%%%%%%%%%%%%%%%%%%%%%%%%%%%%%%%%%%%%%%%%%%%%%%%%%%%%
\begin{frame}\frametitle{Scipy}
    \begin{itemize}
        \item Librería científica
        \item Tiene \emph{muchas} operaciones matemáticas
        \item Álgebra lineal (más eficiente que Numpy)
        \item {\only<2>{\bf}Interpolación de datos}
        \item {\only<2>{\bf}Solución de ecuaciones}
        \item {\only<2>{\bf}Minimización de funciones}
        \item Y MUCHO más
    \end{itemize}

    \begin{flushright}
        \includegraphics[width=0.2\textwidth]{../images/logos/scipy.png}
    \end{flushright}
\end{frame}
%%%%%%%%%%%%%%%%%%%%%%%%%%%%%%%%%%%%%%%%%%%%%%%%%%%%%%%
\begin{frame}\frametitle{Scipy: Interpolación de datos}
    \only<1>{
        \begin{center}
            \includegraphics[width=0.4\textwidth]{../images/interp-data.png}
        \end{center}
        Y si quiero el valor en $x=3$? Eso se llama \emph{interpolación}\footnote{Extrapolación sería bajo 0 o sobre 8}
    }
    \begin{columns}
        \begin{column}{0.6\textwidth}
    \only<2->{
  {\tt\small
  > import numpy as np

  > import matplotlib.pyplot as plt

  > from scipy import interpolate

  > x = np.linspace(0,8,5)

  > y = np.exp(-x/3)

  > f = interpolate.interp1d(x, y, kind)

  > xnew = np.linspace(0,8,100)

  > ynew = f(xnew)
  }
    }
        \end{column}
        \begin{column}{0.39\textwidth}
            \only<2>{
            {\small
            \begin{center}
            \includegraphics[width=\textwidth]{../images/interp-linear}
            \code{kind='linear'}
            \end{center}
            }
            }
            \only<3>{
            {\small
            \begin{center}
            \includegraphics[width=\textwidth]{../images/interp-nearest}
            \code{kind='nearest'}
            \end{center}
            }
            }
            \only<4>{
            {\small
            \begin{center}
            \includegraphics[width=\textwidth]{../images/interp-cubic}
            \code{kind='cubic'}
            \end{center}
            }
            }
        \end{column}
    \end{columns}
\end{frame}
%%%%%%%%%%%%%%%%%%%%%%%%%%%%%%%%%%%%%%%%%%%%%%%%%%%%%%%
\begin{frame}[t,fragile]\frametitle{Scipy: Solución de ecuaciones ($f(x)=0$)}
    \begin{small}
    $$ f(x) \coloneqq x + \sin x  = 0  $$
    \vspace{-0.5cm}
    \begin{columns}
        \begin{column}{0.49\textwidth}
    \begin{minted}{python}  
  from scipy.optimize import root_scalar 
  def f(x): return x+sin(2*pi*x)
  domx = np.linspace(-2,2,1000)
  domy = [f(x) for x in domx]
  x0,x1 = 3, -3
  sol = root_scalar(f, x0=x0, x1=x1)
  plt.scatter(sol.root, f(sol.root))
  plt.plot(domx, domy);plt.show()

    \end{minted}
        \end{column}
        \begin{column}{0.5\textwidth}
            \begin{flushright}
                \includegraphics[width=0.7\textwidth]{../images/scipy-root.png}
            \end{flushright}
        \end{column}
    \end{columns}
    \pause\idea{Cómo obtener las otras soluciones?}
    \end{small}
\end{frame}
%%%%%%%%%%%%%%%%%%%%%%%%%%%%%%%%%%%%%%%%%%%%%%%%%%%%%%%
\begin{frame}[t,fragile]\frametitle{Scipy: Minimización}
    \begin{small}
    $$ \min_ x f(x) \coloneqq \sin^2 4x + x^2  = 0  $$
    \vspace{-0.5cm}
    \begin{columns}
        \begin{column}{0.49\textwidth}
    \begin{minted}{python}  
  from scipy.optimize import minimize_scalar 
  def f(x): return sin(4*x)**2 + x**2
  domx = np.linspace(-4,4,1000)
  domy = [f(x) for x in domx]
  sol = minimize_scalar(f)
  plt.scatter(sol.x, f(sol.x))
  plt.plot(domx, domy);plt.show()
    \end{minted}
        \end{column}
        \begin{column}{0.5\textwidth}
            \begin{flushright}
                \includegraphics[width=0.7\textwidth]{../images/scipy-min.png}
            \end{flushright}
        \end{column}
    \end{columns}
    \end{small}
\end{frame}
%%%%%%%%%%%%%%%%%%%%%%%%%%%%%%%%%%%%%%%%%%%%%%%%%%%%%%%
\darkSlide{Soluciones numéricas NO SON exactas}
%%%%%%%%%%%%%%%%%%%%%%%%%%%%%%%%%%%%%%%%%%%%%%%%%%%%%%%
%%%%%%%%%%%%%%%%%%%%%%%%%%%%%%%%%%%%%%%%%%%%%%%%%%%%%%%
\section{Pandas}
%%%%%%%%%%%%%%%%%%%%%%%%%%%%%%%%%%%%%%%%%%%%%%%%%%%%%%%
%%%%%%%%%%%%%%%%%%%%%%%%%%%%%%%%%%%%%%%%%%%%%%%%%%%%%%%
\begin{frame}[fragile]\frametitle{Tipo central: \code{DataFrame}}
    \begin{minted}{python}
  import pandas
  mydataset = {
    'cars': ["BMW", "Volvo", "Ford"],
    'passings': [3, 7, 2]
  }
  myvar = pandas.DataFrame(mydataset)
  print(myvar) 
    \end{minted}

    \begin{minted}{bash}
  $      cars    passings
    0    BMW         3
    1  Volvo         7
    2   Ford         2
    \end{minted}
\end{frame}
%%%%%%%%%%%%%%%%%%%%%%%%%%%%%%%%%%%%%%%%%%%%%%%%%%%%%%%
\begin{frame}\frametitle{\code{DataFrame}}
    \begin{itemize}
        \item Es el elemento central de Pandas
        \item Son \emph{tablas de datos}
        \item Se pueden filtrar fácilmente
    \end{itemize}
\end{frame}
%%%%%%%%%%%%%%%%%%%%%%%%%%%%%%%%%%%%%%%%%%%%%%%%%%%%%%%
\begin{frame}[fragile]\frametitle{Nombre de datos}

    \begin{minted}{python}
  data = {
    "calories": [420, 380, 390],
    "duration": [50, 40, 45] }

  df = pd.DataFrame(data)
  print(df) 
  df = pd.DataFrame(data, index = ["day1", "day2", "day3"])
  print(df) 
    \end{minted}
\end{frame}
%%%%%%%%%%%%%%%%%%%%%%%%%%%%%%%%%%%%%%%%%%%%%%%%%%%%%%%
\begin{frame}[fragile]\frametitle{Cargar datos}

    Asumiendo que existe archivo \code{data.csv}
    \begin{minted}{python}
  import pandas as pd

  df = pd.read_csv('data.csv')
  print(df)
    \end{minted}

\end{frame}
%%%%%%%%%%%%%%%%%%%%%%%%%%%%%%%%%%%%%%%%%%%%%%%%%%%%%%%
\begin{frame}[fragile]\frametitle{Explorando datos}
    \begin{minted}{python}
  df = pd.read_csv('data.csv')
    \end{minted}
    \begin{itemize}
        \item \code{df.head(n)}: Primeros \code{n} datos
        \item \code{df.tail(n)}: Últimos \code{n} datos
        \item \code{df.describe()}: Ver estadísticas descriptivas
        \item \code{df.sort\_values(by="A")}: Ordenar según columna "A"
    \end{itemize}
    
    \idea{Consola}
\end{frame}
%%%%%%%%%%%%%%%%%%%%%%%%%%%%%%%%%%%%%%%%%%%%%%%%%%%%%%%
\begin{frame}[fragile]\frametitle{Acceso a datos}
    \begin{minted}{python}
  df = pd.read_csv('data.csv')
    \end{minted}
    \begin{itemize}
        \item \code{df["A"]}: Extraer columna con etiqueta "A"
        \item \code{df[1:4]}: Extraer filas en [1,4)
        \item \code{df[1:4][["A", "B"]]}: Filas y columnas
    \end{itemize}

    \idea{Consola}
\end{frame}
%%%%%%%%%%%%%%%%%%%%%%%%%%%%%%%%%%%%%%%%%%%%%%%%%%%%%%%
\begin{frame}[fragile]\frametitle{Filtros booleanos}

    Exploremos los siguientes comandos: 
    \begin{minted}{python}
    df = pd.read_csv('data.csv')
    df["Pulse"] > 120
    df[df["Pulse"] > 120]
    \end{minted}

    \vspace{1cm}
    \pause Se llama filtro por máscara, o máscara booleana.
\end{frame}
%%%%%%%%%%%%%%%%%%%%%%%%%%%%%%%%%%%%%%%%%%%%%%%%%%%%%%%
\begin{frame}\frametitle{Recap}
    \begin{itemize}
        \item Gestión de módulos, PYTHONPATH
        \item \texttt{math, random, statistics, Numpy, Scipy, Matplotlib, Pandas}
        \item Python tiene muchas librerías para hacer ciencia!
        \item También se puede hacer estadística (no lo vimos)
    \end{itemize}

    \pause
    \vspace{1cm}
    \code{https://pandas.pydata.org/docs/user\_guide/10min.html}
\end{frame}
%%%%%%%%%%%%%%%%%%%%%%%%%%%%%%%%%%%%%%%%%%%%%%%%%%%%%%%
\begin{frame}
    \maketitle
\end{frame}
%%%%%%%%%%%%%%%%%%%%%%%%%%%%%%%%%%%%%%%%%%%%%%%%%%%%%%%
\begin{frame}[fragile,noframenumbering]\frametitle{Mini ejercicios}
    \begin{itemize}
        \item 
    \code{operaciones.py}
    \begin{minted}{python}
  def suma(x,y): return x+y

  def mult(x,y): return x*y
    \end{minted}

    \begin{itemize}
        \item Crear el módulo \code{operaciones} e importarlo desde un script.
        \item Guardar el módulo en otro lado e importarlo usando \code{PYTHONPATH}
    \end{itemize}

    \item Genere una sucesión de 0's y 1's que representan lanzamientos de monedas. Calcule indicadores estadísticos de sus datos usando \code{statistics} y \code{pandas}
    \item Considere la función $f(x)=\sin x$ y encuentre numéricamente un mínimo con \code{scipy.minimize}. Qué pasa con los otros mínimos? 
    \end{itemize}
\end{frame}
%%%%%%%%%%%%%%%%%%%%%%%%%%%%%%%%%%%%%%%%%%%%%%%%%%%%%%%

\end{document}
